\documentclass[UTF8,zihao=-4]{ctexart}

\usepackage[a4paper,top=22mm,bottom=22mm,left=22mm,right=22mm]{geometry}
\usepackage{graphicx}
\usepackage{xcolor}
\usepackage{booktabs}
\usepackage{tabularx}
\usepackage{array}
\usepackage{caption}
\usepackage{subcaption}
\usepackage{float}
\usepackage[section]{placeins}
\usepackage{pdflscape}
\usepackage{enumitem}
\usepackage{fancyhdr}
\usepackage{hyperref}

\definecolor{reportblue}{HTML}{365F91}
\definecolor{reportgreen}{HTML}{5E8C61}
\definecolor{reportorange}{HTML}{D9822B}
\definecolor{reportred}{HTML}{B64B4B}
\definecolor{reportgray}{HTML}{5F6368}
\definecolor{reportlight}{HTML}{F3F5F7}

\hypersetup{
  colorlinks=true,
  linkcolor=reportblue,
  urlcolor=reportblue,
  pdfauthor={韦承谦},
  pdftitle={藕塘滑坡位移概率预测与多测点预警研究进展报告}
}

\captionsetup{font=small,labelfont=bf,labelsep=quad}
\setlength{\parindent}{2em}
\setlength{\parskip}{0.35em}
\setlist[itemize]{leftmargin=2em,itemsep=0.25em,topsep=0.35em}
\renewcommand{\arraystretch}{1.25}

\pagestyle{fancy}
\fancyhf{}
\fancyhead[L]{藕塘滑坡研究进展报告}
\fancyhead[R]{阶段稿}
\fancyfoot[C]{\thepage}
\setlength{\headheight}{15pt}

\newcommand{\statusbox}[2]{%
  \noindent\colorbox{reportlight}{%
    \parbox{0.96\linewidth}{\textbf{#1}\quad #2}%
  }%
}

\title{\heiti\LARGE 藕塘滑坡位移概率预测与多测点预警研究进展报告\\[0.6em]
\large 当前工程原型、结果判断与待确认问题}
\author{韦承谦}
\date{2026年8月}

\begin{document}

\maketitle

\statusbox{报告定位}{本报告用于说明藕塘案例目前已经完成什么、结果是否达到预期，以及正式方法还缺少什么。当前结果是可复算的工程原型，不是现场正式预警。}

\vspace{0.8em}

\statusbox{数据边界}{现有输入是公开发布的物化日序列。原始 GNSS 锚点和日值生成链无法取得，因此可用于流程跑通和内部诊断，但不能作为独立原始监测数据上的确认性证据。}

\vspace{0.8em}

\statusbox{案例范围}{本阶段只使用藕塘 8 个测点。Vajont 尚未启动；后续读取、适配或运行仍须获得明确授权。}

\clearpage

\section{本阶段要回答的问题}

\emph{本节将在下一提交中加入技术路线总览图和简要说明。}

\section{ConvLSTM 概率位移预测}

\emph{本节将在下一提交中加入全部测点预测图和三折×五种子结果图。}

\section{SHAP 模型依赖分析}

\emph{本节将在下一提交中加入 SHAP 稳定性和删组结果。}

\section{四指标与多测点预警原型}

\emph{本节将在下一提交中加入完整时间线、逐点联合诊断和典型日空间融合图。}

\section{当前结论与下一步}

\emph{本节将在下一提交中区分已完成、未完成和需要导师确认的事项。}

\end{document}
